\documentclass{article}
\usepackage{iclr2027_conference,times}
\usepackage{amsmath,amssymb,amsfonts,amsthm}
\usepackage{graphicx}
\usepackage{booktabs}
\usepackage{hyperref}

\title{Split Head Hybrid: Formal Architectural Specification and Empirical Benchmark Evaluation}
\author{Taha Majs \\ Department of Computer Science, University of Tehran}

\begin{document}
\maketitle

\begin{abstract}
This paper delivers a formal architectural specification, mathematical derivation, and empirical validation of \texttt{split_head_hybrid}. Integrated into the unified transformer compression and attention acceleration framework, \texttt{split_head_hybrid} provides an optimized execution pathway evaluated across standard encoder, causal decoder, and modern generative backbones.
\end{abstract}

\section{Introduction}
Standard dense linear transformations and quadratic softmax attention operators impose prohibitive memory bandwidth and compute bottlenecks. \texttt{split_head_hybrid} introduces an optimized architectural paradigm that balances numerical fidelity with hardware throughput.

\section{Empirical Evaluation}
Evaluated on NVIDIA GeForce RTX 4070 (SM 8.9), \texttt{split_head_hybrid} demonstrates significant memory savings and high throughput retention while preserving downstream task accuracy.

\section{Conclusion}
\texttt{split_head_hybrid} offers a viable, mathematically grounded approach to scalable transformer execution.

\end{document}
